\documentclass[11pt]{article}
\usepackage[margin=1in]{geometry}
\usepackage[T1]{fontenc}
\usepackage{xcolor}
\usepackage{amsmath,amssymb,booktabs,tabularx,enumitem,array,hyperref,microtype}
\hypersetup{colorlinks=true,linkcolor=blue!50!black}
\newcolumntype{L}{>{\raggedright\arraybackslash}X}
\setlength{\parindent}{0pt}\setlength{\parskip}{6pt}
\title{Positioning Embeddings, Attention and LightGBM:\\ Methodology}
\author{PositioningWork project}
\date{1 November 2026}
\begin{document}
\maketitle

\textbf{Status: a design only. Nothing here is built or tested.} All sizes and settings are proposals. The design has no published precedent that we found. The evidence so far is weak: no positioning signal has passed an out-of-sample test in our studies. So every step must earn its place against a simpler model.

\section*{Summary}
\begin{enumerate}
\item Turn positioning into features, keeping the full shape of the ledger, not one Pain number.
\item Test those features with ridge regression.
\item Test them with a small LightGBM (LGBM) model.
\item Add one small attention layer, freeze it, and pass its outputs to LGBM.
\item Keep a step only if it beats the step before it, out of sample.
\end{enumerate}

\textbf{Why attention at all.} The literature suggests that a loss matters when it meets a binding limit, and that holder type can change the sign (Siriwardane 2019; Ellul et al.\ 2015; Kang et al.\ 2020; Fang 2024). Attention lets market state decide which reading to trust. LGBM can also learn such interactions, so attention may add nothing. Step 4 tests that.

\section{Notation and data}
A \emph{unit} $u$ is what we forecast, such as a bond, an issuer, a sector and maturity bucket, an ETF or an index. Time $t$ is weekly. Every input must be known at $t$. Each feature carries the date it became available.

\begin{center}\small
\begin{tabularx}{\linewidth}{@{}lLL@{}}
\toprule
Market & Available positioning inputs & Limit on this method \\
\midrule
TRACE & Customer buys and sells by bond & One holder group only, so attention can only weigh measures. Our pilot panel has 12 issuers and 123 weeks, and it is not point-in-time. \\
RFQ (work machine) & Anonymous side, size and level, from one dealer & No client view. Use buckets. Any signal must beat a vendor mark. \\
Credit ETFs & Shares, NAV, FINRA short interest & Share count is not a position. Prefer baskets and premiums if available. \\
Credit futures & CFTC holder groups & About 25 reports. Monitor only. Do not train. \\
CDX & Prices, volumes, skew, cleared totals & No public holder data. Proxies only. \\
\bottomrule
\end{tabularx}
\end{center}

\section{Step 1: features}
\subsection{The ledger-shape embedding}
This is the main new feature. For holder group $g$ and unit $u$, build the lot ledger from net flow, with pro-rata removal. Entry level is the trade's spread, or total-return price for an ETF. For each open long lot $\ell$, compute its distance from today, in volatility units:
\[
D_\ell = \frac{s_t - s_\ell}{\sigma_t} \quad\text{for spreads}, \qquad
D_\ell = \frac{P_\ell - P_t}{\sigma_t} \quad\text{for prices}.
\]
A positive distance means the lot is under water. Short lots use the opposite sign. $\sigma_t$ is the standard deviation of weekly changes over the past 26 weeks.

Place the lots in 16 bins, with edges from $-3.5$ to $+3.5$ in steps of $0.5$, plus two open ends. Each bin holds that bin's share of open quantity:
\[
h_{b} = \frac{\sum_{\ell \in b} q_\ell}{\sum_{\ell} q_\ell}, \qquad \sum_b h_b = 1 .
\]
Pain share and depth are summaries of $\mathbf{h}$. So $\mathbf{h}$ keeps everything they hold, and more.

\textbf{Rules.}
\begin{itemize}
\item If open quantity is $10^{-6}$ or less, mark the row missing. Add a missing flag.
\item Always add coverage, which is open quantity over the amount outstanding or shares.
\item Starting positions are unknown, so read $\mathbf{h}$ as the shape of recent flow, not of all holders.
\item Run one test with Grinblatt and Han (2005) turnover decay, which also retires lots with volume. Our current ledger does not do this.
\end{itemize}

\subsection{Other readings}
Build these measures per group where the data allows.
\begin{center}\small
\begin{tabular}{@{}ll@{}}
\toprule
Measure & Source in the spec \\
\midrule
Net position scaled & I1, I3 \\
Range position & I2 sentiment index, if history allows \\
Flow trend & TAARSS from I13 \\
Filtered residual & I5 \\
Directional short interest & ETFs only, following Evans et al.\ (2026) \\
\bottomrule
\end{tabular}
\end{center}
Standardise each reading $x$ with an expanding mean and standard deviation, known at $t$. Clip it at $\pm 3$.

\subsection{Market state}
The state vector $\mathbf{s}_t$ has five inputs, all standardised: spread level, spread volatility, past 26-week spread change, a dealer capacity proxy, and a funding stress proxy. Choose and freeze the two proxies before any test.

\subsection{Target}
The target is the future spread change over $H$ weeks, for $H \in \{1, 4, 13\}$, less the change in the unit's market bucket. Start with $H=4$. Use spreads or excess returns, not price, because rates and pull-to-par distort price. Follow Dickerson et al.\ on bond data: avoid one-sided filters and allow for price error.

\section{Step 2: ridge}
Fit ridge regression on $\mathbf{h}$ for each group, the other readings, and $\mathbf{s}_t$. Choose the penalty inside the training data. This model tells you whether any linear signal exists. If it does not beat the benchmarks in Section~\ref{sec:eval}, record that, and treat steps 3 and 4 as exploratory.

\section{Step 3: LightGBM}
Use the same inputs. Add hand-built interactions, such as depth $\times$ dealer capacity, and depth by holder group. This is the simple alternative to attention, and it has far less leakage risk.

\begin{center}\small
\begin{tabular}{@{}ll@{}}
\toprule
Setting & Proposed value \\
\midrule
Maximum depth, leaves & 3, up to 8 \\
Minimum rows per leaf & 5\% of training rows \\
Learning rate & 0.02, with early stopping on a later block \\
Feature fraction & 0.5 \\
L2 penalty & 10 \\
Row weights & Each week sums to 1, so busy weeks do not dominate \\
Splits & By whole weeks, never by row \\
\bottomrule
\end{tabular}
\end{center}

\section{Step 4: attention}
\subsection{The layer}
Each token $i$ is one reading, with holder group $g_i$ and measure $m_i$. Its tag, in $d = 8$ dimensions, is learned:
\[
\mathbf{e}_i = \mathbf{e}^{G}_{g_i} + \mathbf{e}^{M}_{m_i}.
\]
The query comes from market state:
\[
\mathbf{q}_t = \mathbf{W}\mathbf{s}_t + \mathbf{c}.
\]
The attention weights and the output are:
\[
a_{i,t} = \frac{\exp(\mathbf{q}_t^\top \mathbf{e}_i / \tau)}{\sum_j \exp(\mathbf{q}_t^\top \mathbf{e}_j / \tau)},
\qquad
z_{u,t} = \sum_i a_{i,t}\, b_{g_i,m_i}\, x_{i,u,t} .
\]
The ledger shape enters as a single reading per group, $x = \mathbf{w}^\top \mathbf{h}$, with one vector $\mathbf{w}$ shared across groups. Set $\tau = 1$.

\textbf{Why the signed loadings $b$.} Weights $a$ are always positive. Without $b$, the layer could not let insurers and levered holders carry opposite signs.

\subsection{Training}
A small linear head forecasts the target:
\[
\hat{y}_{u,t} = \beta_0 + \beta z_{u,t} + \boldsymbol{\gamma}^\top \mathbf{s}_t .
\]
Minimise Huber loss, plus an L2 penalty on all parameters, plus a small penalty for over-concentrated weights. Use Adam, whole-week batches and early stopping on a later block. Use three random seeds, and average the frozen outputs.

\subsection{Size limit}
With 3 groups, 5 measures, 16 bins and 5 state inputs, the layer has about 150 parameters:
\begin{center}\small
\begin{tabular}{@{}lr@{}}
\toprule
Part & Parameters \\
\midrule
Tags, $(3+5)\times 8$ & 64 \\
Query, $5\times 8 + 8$ & 48 \\
Signed loadings, $3\times 5$ & 15 \\
Ledger-shape vector & 16 \\
Head, $1 + 1 + 5$ & 7 \\
\midrule
Total & 150 \\
\bottomrule
\end{tabular}
\end{center}
Do not exceed this without new evidence. Use one head and one layer.

\subsection{Passing it to LGBM}
LGBM cannot train the layer, because trees pass no gradients back. So freeze the layer, then export these columns:
\begin{enumerate}
\item The weights $a_{i,t}$.
\item The output $z_{u,t}$.
\item The weight entropy $-\sum_i a_{i,t}\ln a_{i,t}$, which shows whether one reading dominates.
\end{enumerate}
Add these to the inputs from Step 3. Do not export $\hat{y}$ as well, as it adds little beyond $z$ and $\mathbf{s}_t$.

\section{Preventing leakage}
This is the main risk. If attention trains on the weeks LGBM learns from, its columns carry those weeks' targets, and LGBM will look far better than it is.

\begin{enumerate}
\item Set a final holdout before any fit. No model or choice touches it until the end.
\item Use an expanding walk-forward, refitting yearly. Everything fits only on data before the cut.
\item Inside each training window, split the weeks into five contiguous blocks.
\item For each block, fit attention on the other blocks, and write its columns for that block. Remove $H$ weeks on each side of the block first. These out-of-fold columns are LGBM's training inputs.
\item Then fit attention on the whole training window. Use that fit to write the columns for test weeks.
\item Leave a gap of $H$ weeks between training data and test data.
\item Compute all standardising, binning and volatility on past data only.
\item TRACE in our warehouse is not point-in-time. Label results from it as provisional.
\end{enumerate}

\section{Evaluation}\label{sec:eval}
\subsection{The ladder}
\begin{center}\small
\begin{tabular}{@{}lll@{}}
\toprule
Model & Inputs & Question \\
\midrule
B0 & Spread level and past spread change & The bar to beat \\
M1 & Ridge on all features & Is there linear signal? \\
M2 & LGBM on the same features, plus hand-built interactions & Do interactions help? \\
M3 & M2 plus frozen attention columns & Does attention add anything? \\
\bottomrule
\end{tabular}
\end{center}
Each model must beat the one before it. The plain spread level had $t=5.93$ on LQD and $12.51$ on HYG at 26 weeks in our ETF study, so B0 is a demanding bar.

\subsection{Statistics}
\begin{itemize}
\item Primary measure: weekly cross-sectional rank correlation with the target, averaged over time. Use Newey-West errors with lag at least $H-1$.
\item Few issuers means weak clustering. Use a wild cluster bootstrap by issuer and date.
\item Report effective independent observations for each horizon.
\item Placebos: circularly shift flow within each unit, and shuffle holder labels. M3 must beat both.
\item Apply a Holm correction across all models, horizons and markets. Report raw and adjusted p-values.
\item Report a net result, after stated costs, for any trading claim.
\end{itemize}

\subsection{Checks that attention means something}
\begin{enumerate}
\item Weights should be stable across seeds and yearly refits. If they are not, do not read them.
\item Check whether weights shift sensibly in known stress, without using that check to tune the model.
\item Compare SHAP importance for attention columns against the placebo versions.
\end{enumerate}

\subsection{Stop rules}
\begin{enumerate}
\item If M1 and M2 fail against B0, stop. Attention cannot create information that the inputs lack.
\item If M3 does not beat M2 out of sample, drop attention and keep the simpler model.
\item If weights are unstable, drop attention even if M3 scores slightly better.
\item Do not widen the search after seeing the holdout.
\end{enumerate}

\section{Order of work}
\begin{enumerate}
\item The TRACE pilot passed 0 of 48 tests on 12 large issuers, and flow looked like noise. So do not start on that panel. Start only on a much larger one.
\item Build $\mathbf{h}$ on TRACE, then run B0 and M1.
\item If M1 passes, run M2.
\item If M2 passes, run M3 on a larger sample. Our current TRACE panel is too small to train attention reliably, with 123 weeks and 12 issuers. Treat any M3 result on it as exploratory.
\item Repeat the ladder on work-machine RFQ data, using buckets.
\end{enumerate}

\section*{Sources}
Full citations and reading levels are in \texttt{../literature/literature\_review.pdf}. The ledger, measures and their recipes are in \texttt{../../positioning\_spec.pdf}.
\end{document}
